%% FrankenMSA -- Supplementary Information
%% Compile BEFORE main_oup.tex: main_oup.tex reads the labels of this document
%% (via the xr package, prefix "S-") from supplementary.aux.
%%
%% Markers:  \TODO{...}       red   -- missing information to be filled in
%%           \citeNEW{...}    blue  -- citation newly added in this revision (check!)

\documentclass[11pt,a4paper]{article}

\usepackage[margin=2cm]{geometry}
\usepackage[T1]{fontenc}
\usepackage{lmodern}
\usepackage{graphicx}
\usepackage{subfig}
\usepackage{booktabs}
\usepackage{amsmath}
\usepackage{xcolor}
\usepackage{url}
\usepackage[hidelinks]{hyperref}

\newcommand{\TODO}[1]{\textcolor{red}{[TODO: #1]}}

% Supplementary numbering: Section S1, Fig. S1, Table S1
\renewcommand{\thesection}{S\arabic{section}}
\renewcommand{\thesubsection}{\thesection.\arabic{subsection}}
\renewcommand{\thefigure}{S\arabic{figure}}
\renewcommand{\thetable}{S\arabic{table}}
\renewcommand{\figurename}{Supplementary Fig.}
\renewcommand{\tablename}{Supplementary Table}

\title{Supplementary Information for\\[0.3em]
\textbf{FrankenMSA: an interactive toolbox for engineering multiple sequence alignments}}
\author{\TODO{authors}}
\date{}

\begin{document}
\maketitle

\section{Software}\label{sec:install}
\subsection{Installation}
The modules for ESM3-based clustering \cite{esm3} and ProteinMPNN-based inverse folding \cite{proteinmpnn} depend on large neural networks; they run substantially faster on a GPU and can be deselected during installation. Homology search with MMseqs2 uses the public ColabFold server \cite{colabfold}, and PLMSearch \cite{plmsearch} uses their server; both therefore require an internet connection and send the query sequences to the respective server.
Visit the GitHub page to view the installation instructions and to run the web application locally or in Google Colab: \url{https://github.com/ibmm-unibe-ch/FrankenMSA/}.
\begin{figure}[p]
\centering
\subfloat[Start page]{\includegraphics[width=0.77\textwidth]{img/screenshot_landing.png}}\\[0.5em]
\subfloat[Edit module]{\includegraphics[width=0.77\textwidth]{img/screenshot_edit.png}}\\[0.5em]
\subfloat[Align module]{\includegraphics[width=0.77\textwidth]{img/screenshot_align.png}}
\caption{Graphical user interface of FrankenMSA. \textbf{a} The start page; the navigation bar at the top gives access to the eight modules, and the drop-down menu on the right selects the active MSA and shows its dimensions. \textbf{b} The Edit module summarises the selected MSA (number of sequences, sequence lengths, gaps, consensus and query sequence) and provides tools to filter, sort and shuffle, slice and crop, and edit it. \textbf{c} The Align module generates MSAs with the ColabFold MMseqs2 server, for monomers or as paired MSAs for multimers, or with PLMSearch.}
\label{fig:gui}
\end{figure}


\section{Alternative conformations}\label{sec:altconf}
Both engineered MSAs combine sequences from the ColabFold MMseqs2 server \cite{mmseqs2,colabfold} with synthetic sequences designed by ProteinMPNN \cite{proteinmpnn} for the target conformation \cite{inverse_folding}.
The two sources are joined position-wise: for each region of the query, all rows of the MSA are taken from either the MMseqs2 MSA or the synthetic MSA (bars in Fig.~1 of the main text). We generated 128 sequences with ProteinMPNN at sampling temperature 1.0 for each target conformation and subsampled the MMseqs2 MSAs to contain 128 sequences. The construction of both MSAs is described in the following subsections.

\subsection{Pheromone-binding protein of \textit{Bombyx mori}}
The MMseqs2 MSA was generated for the full sequence of chain A of PDB entry 2FJY. The synthetic MSA consisted of 128 sequences designed by ProteinMPNN for the full structure of 2FJY chain A.
The engineered MSA was taken from the synthetic MSA for residues 0--16, 55--83 and 123--133 and from the MMseqs2 MSA for all remaining residues.
Structures were predicted with AlphaFold2 \cite{alphafold2}.

\subsection{Fusion glycoprotein F0 of respiratory syncytial virus}
The query comprised residues 25--510 of chain A of PDB entry 5U68, i.e.\ the full sequence without the parts that are not modelled in the structure, followed by the last 64 residues of UniProt entry P03420, which form the transmembrane domain. The MMseqs2 MSA was generated for this sequence and triplicated to form the trimer.
The synthetic MSA was generated by inverse folding chain A of 5U68 with ProteinMPNN. The engineered MSA was taken from the synthetic MSA for residues 0--100 and 135--486 and from the MMseqs2 MSA for residues 100--135, a region that is not resolved in 5U68, and for the transmembrane domain (residues 486--550).
Structures were predicted with AlphaFold3 \cite{alphafold3}.

\subsection{Evaluation}
\TODO{maybe how predictions were compared with the experimental conformations, e.g.\ TM-score and RMSD of the top-ranked model to 2FJY/2P71 and 5U68/3RRR; add a table with the values}

\section{Effects of point mutations}\label{sec:mutations}
\subsection{Dataset}
We used the S669 dataset of 669 single-point mutations with experimentally determined stability changes \cite{stability_dataset}, in the preprocessed version of \cite{stability_dataset_fixed}.
Throughout, positive $\Delta\Delta G$ denotes stabilising and negative $\Delta\Delta G$ destabilising mutations.

For each of the 669 mutations and wild-type sequences, the full MSA was generated with the ColabFold MMseqs2 server \cite{mmseqs2,colabfold}.
The narrow MSA was obtained by filtering the full MSA for sequences that carry either the wild-type or the mutant residue at the mutated position.
Wild type and mutant were predicted with AlphaFold2 \cite{alphafold2} from the full MSA, from the narrow MSA and from the query sequence alone, each with three recycles, one seed, and without templates.
For each input, we extracted the local pLDDT at the mutated position from the top-ranked model and computed $\Delta$local pLDDT as the value of the mutant minus that of the wild type.

As baselines, we used the structure-based stability predictor MAESTRO \cite{maestro} and the empirical force field FoldX \cite{FoldX_1,FoldX_2} adjacent to the dataset \cite{stability_dataset}.
Their predicted $\Delta\Delta G$ values are shown against the experimental values in \autoref{fig:competitors}.

Mutations with $\Delta\Delta G \geq +1$~kcal\,mol$^{-1}$ were classified as stabilising (positive class) and mutations with $\Delta\Delta G \leq -1$~kcal\,mol$^{-1}$ as destabilising (negative class); mutations with $|\Delta\Delta G| < 1$~kcal\,mol$^{-1}$ were excluded from the classification to account for the noise in the experimental data.
A prediction counts as positive if its score is greater than zero, where the score is $\Delta$local pLDDT for the AlphaFold2 inputs and the predicted $\Delta\Delta G$ for MAESTRO and FoldX.
\TODO{Check what to report exactly} Precision, recall and F1 refer to the positive class, specificity and negative predictive value (NPV) to the negative class, and the macro-averaged F1 is the mean of the F1 scores of both classes. AUC is the area under the receiver operating characteristic curve of the raw score on the classified mutations, and $\rho_\mathrm{all}$ is the Spearman correlation between score and experimental $\Delta\Delta G$ over all mutations. The results are listed in Table~\ref{tab:ddg_classification}.

\begin{figure}[t]
\centering
\subfloat[MAESTRO]{\includegraphics[width=0.45\textwidth]{img/scatter_maestro.png}}\hfill
\subfloat[FoldX]{\includegraphics[width=0.45\textwidth]{img/scatter_foldx.png}}
\caption{Predicted against experimental stability change $\Delta\Delta G$ for the S669 dataset using \textbf{a} MAESTRO and \textbf{b} FoldX. Positive $\Delta\Delta G$ denotes stabilising mutations.}
\label{fig:competitors}
\end{figure}

% Table generated by the analysis script; regenerate img/classification_metrics.tex
% after the update to all 669 mutations (also update the caption text there).
% It defines \label{tab:ddg_classification}.
\input{img/classification_metrics.tex}

\clearpage
\bibliographystyle{oup-plain}
\bibliography{biblio}

\end{document}
